\section{Introduction}
\label{sec:introduction}

Low Earth orbit (LEO) satellites extend Internet-of-Things (IoT)
connectivity to areas beyond the coverage of terrestrial networks.
The sporadic transmission of short packets by a large population of
terminals makes grant-free random access (GFRA) attractive for satellite
IoT, since active terminals can transmit without first obtaining an
uplink grant \cite{shen2025}. With limited pilot resources, the receiver
must identify active terminals and estimate their channels from
superimposed nonorthogonal pilot signals. This gives rise to the joint
activity detection and channel estimation (JADCE) problem
\cite{liu2018}. In LEO uplinks, the satellite motion and the different
propagation distances across its coverage area introduce user-dependent
delays and Doppler shifts, which further complicate pilot reception.

The Release~17 New Radio non-terrestrial network (NR-NTN) framework
assumes terminals with global navigation satellite system (GNSS)
capabilities. A terminal uses its GNSS position, satellite ephemeris,
and common timing advance information to perform uplink timing and
frequency pre-compensation \cite[Secs.~16.14.1 and 16.14.2.2]{3gpp38300}.
For low-cost, battery-powered IoT terminals, GNSS acquisition and the
associated compensation procedures impose additional hardware and
energy requirements \cite{shen2025}. This motivates GFRA schemes that
operate without terminal GNSS assistance. In the setting considered
here, a common reference delay is compensated, while the satellite
receiver handles the remaining differential delays and Doppler shifts.
These offsets can exceed a symbol duration and a subcarrier spacing,
respectively, and therefore need to be retained in the signal model
\cite{shen2025}.

Waveforms designed for doubly dispersive channels have been investigated
for satellite random access. Shen et al.\ proposed a spreading-based
multi-frame orthogonal time-frequency space (OTFS) transmission scheme
and a receiver for joint device identification, channel estimation, and
symbol detection in the presence of large differential delays and
Doppler shifts \cite{shen2025}. Affine frequency division multiplexing
(AFDM) provides another approach through the discrete affine Fourier
transform (DAFT), whose chirp parameters can be selected according to
the delay and Doppler spreads \cite{bemani2022}. For AFDM-based LEO
GFRA, Ding et al.\ developed a graph neural network that refines sparse
priors by learning correlations among delay--Doppler grid nodes
for JADCE \cite{ding2026}.

The choice of channel representation determines which dependencies a
receiver can exploit. A discrete delay--Doppler representation describes
the channel through path coefficients on a prescribed grid. Fractional
delays and Doppler shifts spread the contribution of a physical path
over multiple grid points, making the assumed support less accurate
\cite{shen2025,huang2026}. An alternative is to estimate the AFDM
effective channel directly \cite{huang2026}. In a multiantenna system,
each terminal's effective channel comprises a collection of matrices
that share path delays, Doppler shifts, and array responses. These
physical parameters induce dependencies among the matrix entries and
across receive antennas. Learning their joint distribution offers a way
to use these dependencies in channel estimation without prescribing a
discrete path support.

Score-based generative models learn a channel prior through the gradient
of the log density of Gaussian-perturbed channel samples. Arvinte and
Tamir applied such models to multiple-input multiple-output (MIMO)
channel estimation using posterior sampling \cite{arvinte2023}.
For LEO AFDM links, Huang et al.\ learned the score of the effective
channel and performed inference using likelihood-annealed Langevin
dynamics \cite{huang2026}. To incorporate generative priors into JADCE,
Cai et al.\ combined channel denoising with message passing (MP) in a
MIMO orthogonal frequency division multiplexing (OFDM) system
\cite{cai2026}. Their receiver uses first- and second-order score
networks to obtain channel moment estimates at each iteration, avoiding
an online reverse diffusion sampling procedure.

In the GNSS-independent AFDM uplink, the observation model must also
account for the unknown phase increments accumulated between pilot
frames. These increments enter the sensing matrix, whereas the
within-frame delay and Doppler effects enter the effective channel.
Joint recovery of terminal activity and channels also requires evidence
for the active and inactive hypotheses. We combine the learned
channel moments with a density ratio for activity detection, allowing
the MP receiver to use the joint distribution of complete multiantenna
AFDM channel blocks for both channel estimation and activity detection.

In this paper, we investigate generative-model-aided GFRA for a
GNSS-independent multiuser MIMO-AFDM satellite IoT system, where
single-antenna terminals communicate with a LEO satellite equipped with
a uniform planar array (UPA). Our main contributions are as follows.
\begin{enumerate}[label=\arabic*)]
\item We formulate a multi-frame AFDM uplink model with continuous path
delays and Doppler shifts. A chirp-periodic prefix (CPP) accommodates
the differential propagation delay, and the model separates
inter-frame phase evolution from the reference-frame effective
channel. Stacking the received pilots yields a multiuser observation
model for joint activity detection and full multiantenna channel
estimation.

\item We develop a generative-model-aided MP receiver for JADCE.
The linear estimator processes the multiuser pilot observations, while
the nonlinear denoiser combines the learned channel prior with the
terminal activity model. The two modules iteratively exchange
extrinsic means and variances to update the channel estimates and
activity probabilities. The inter-frame phase increments and noise
variance are estimated from the received pilots before the MP iterations.

\item We incorporate a learned prior for the complete multiantenna
AFDM effective channel into the nonlinear denoiser. Two score networks
provide a channel mean estimate and a scalar error estimate. A
separately trained classifier estimates the density ratio between
active and inactive observations \cite{menon2016}, which supplies the
activity probabilities used in the channel estimates and outgoing
messages. We specify the training objectives, message schedule, and
computational complexity of the resulting receiver.
\end{enumerate}

Numerical comparisons over 500 independent physical scenes show lower
channel estimation error than the four evaluated baseline receivers
over the tested signal-to-noise ratio (SNR) range of $-4$ to $8$~dB,
and lower activity detection error over $-2$ to $8$~dB.

The remainder of this paper is organized as follows.
Section~\ref{sec:system} presents the system model and problem
formulation. Section~\ref{sec:algorithm} develops the receiver, and
Section~\ref{sec:learning} describes the learning objectives.
Section~\ref{sec:overall} summarizes the algorithm and analyzes its
complexity. Numerical results and conclusions are presented in
Sections~VI and VII, respectively.
